\centerline{\bf \Huge Можно ли? II}

\begin{flushright}
        \scriptsize{Можно все, что хочешь, когда хочешь только то, что следует.}
\end{flushright}

\problem{\textbf{1}}  На доске написано число 49. За один ход разрешается либо удваивать число, либо стирать его последнюю цифру. Можно ли за несколько ходов получить число 50?

\problem{\textbf{2}} Можно ли в какой-то прямоугольной таблице расставить натуральные числа так, чтобы в каждом столбце сумма чисел была больше 150, а в каждой строке – меньше 8?

\problem{\textbf{3}} Можно ли разрезать квадрат на квадратики двух размеров так, чтобы маленьких было столько же, сколько и больших?

\problem{\textbf{4}} Алина разлила сок на клетчатый лист тетради размером 100 × 101 клеток.
Могло ли после этого получиться, что испачканных клеток на 2025 больше, чем чистых?


\problem{\textbf{5}} Имеется 12 сосисок длиной 13 см каждая. Их
требуется разделить между 13 котятами, 13 кошками и 13 котами так, чтобы каждому
котенку достался кусок сосиски длиной 3 см, каждой кошке — кусок сосиски длиной
4 см, а каждому коту кусок сосиски длиной 5 см. Можно ли это сделать?


\problem{\textbf{6}} Начальник выписал на доску в ряд некоторые целые числа от 0 до
999. В итоге получилось длинное число. НеНачальник записал на свою часть доски все оставшиеся
целые числа из этого же диапазона, в итоге получилось второе длинное число. Могли ли эти
два длинных числа совпасть?


\problem{\textbf{7}} Магический квадрат — это квадрат, составленный из чисел, в котором суммы чисел в каждой
строке, в каждом столбце и на каждой диагонали равны. Можно ли составить магический
квадрат 5 $\times$ 5 из первых 25 простых чисел?

